%% arara directives
% arara: xelatex
% arara: bibtex
% arara: xelatex
% arara: xelatex

%\documentclass{article} % One-column default
\documentclass[twocolumn, switch]{article} % Method A for two-column formatting
\usepackage{float}
\usepackage{graphicx}
\usepackage{subcaption}
%% Math packages
\usepackage{amsmath, amsthm, amssymb, amsfonts}
\usepackage{bigints}
%% Bibliography options
\usepackage[numbers,square]{natbib}
\bibliographystyle{unsrtnat}
%\usepackage{natbib}
%\bibliographystyle{Geology}
\usepackage{mwe}
\usepackage{appendix}
%% General packages
\usepackage[utf8]{inputenc}	% allow utf-8 input
\usepackage[T1]{fontenc}	% use 8-bit T1 fonts
\usepackage{xcolor}		% colors for hyperlinks
\usepackage[colorlinks = true,
            linkcolor = purple,
            urlcolor  = blue,
            citecolor = cyan,
            anchorcolor = black]{hyperref}	% Color links to references, figures, etc.
\usepackage{booktabs} 		% professional-quality tables
\usepackage{nicefrac}		% compact symbols for 1/2, etc.
\usepackage{microtype}		% microtypography
\usepackage{lineno}		% Line numbers
\usepackage{float}			% Allows for figures within multicol
%\usepackage{multicol}		% Multiple columns (Method B)

\usepackage{lipsum}		%  Filler text
\usepackage{preprint}
 %% Special figure caption options
\usepackage{newfloat}
\DeclareFloatingEnvironment[name={Supplementary Figure}]{suppfigure}
\usepackage{sidecap}
\sidecaptionvpos{figure}{c}



% Section title spacing  options
\usepackage{titlesec}
\titlespacing\section{0pt}{12pt plus 3pt minus 3pt}{1pt plus 1pt minus 1pt}
\titlespacing\subsection{0pt}{10pt plus 3pt minus 3pt}{1pt plus 1pt minus 1pt}
\titlespacing\subsubsection{0pt}{8pt plus 3pt minus 3pt}{1pt plus 1pt minus 1pt}

% ORCiD insertion
\usepackage{tikz,xcolor,hyperref}

\definecolor{lime}{HTML}{A6CE39}
\DeclareRobustCommand{\orcidicon}{
	\begin{tikzpicture}
	\draw[lime, fill=lime] (0,0) 
	circle [radius=0.16] 
	node[white] {{\fontfamily{qag}\selectfont \tiny ID}};
	\draw[white, fill=white] (-0.0625,0.095) 
	circle [radius=0.007];
	\end{tikzpicture}
	\hspace{-2mm}
}
\foreach \x in {A, ..., Z}{\expandafter\xdef\csname orcid\x\endcsname{\noexpand\href{https://orcid.org/\csname orcidauthor\x\endcsname}
			{\noexpand\orcidicon}}
}
% Define the ORCID iD command for each author separately. Here done for two authors.
\newcommand{\orcidauthorA}{0009-0004-8185-5891}
\newcommand{\orcidauthorB}{0000-0002-6168-9142}
\newcommand{\orcidauthorC}{0000-0000-0000-0003}
\newcommand{\orcidauthorD}{0000-0000-0000-0004}

%%%%%%%%%%%%%%%%   Title   %%%%%%%%%%%%%%%%
\title{Unconventional optical response in monolayer graphene upon dominant intraband scattering}

%%%%%%%%%%%%%%%  Author list  %%%%%%%%%%%%%%%
\usepackage{authblk}
\renewcommand*{\Authfont}{\bfseries}
\author[]{Palash Saha\orcidA{}}
\author[]{Bala Murali Krishna Mariserla*\orcidB{}}


\affil[]{Ultrafast Physics Group, Department Of Physics, Indian Institute Of Technology Jodhpur, 342037, India.}

\affil{*bmkrishna@iitj.ac.in}

%%%%%%%%%%%%%%    Front matter    %%%%%%%%%%%%%%
\begin{document}

\twocolumn[ % Method A for two-column formatting
  \begin{@twocolumnfalse} % Method A for two-column formatting
  
\maketitle

\begin{abstract}
Scattering kinetics influence the graphene's transport properties and inhibits the charge carrier deterministic behaviour. The intra/inter-band scattering mechanisms are vital for graphene’s optical conductivity response under specific considerations of doping. Here, we investigated the influence of scattering  systematically on optical conductivity using the multiband semi-classical Boltzmann equation with inclusion of both electron-electron $\&$ electron-phonon collisions. We found unconventional characteristics of linear optical response  with a significant deviation from the universal conductivity $(e^2/4\hbar)$ in doped monolayer graphene. This is examined through phenomenological relaxation rates under low doping regimes and found that the dominance of intraband scattering . Such novel optical responses are vanished at high temperatures or overdoping conditions due to strong Drude behaviour. With the aid of approximations around Dirac points we have developed formalism for many body interactions and found which is in good agreement with the Kubo approaches.
\end{abstract}
%\keywords{First keyword \and Second keyword \and More} % (optional)
\vspace{0.35cm}

  \end{@twocolumnfalse} % Method A for two-column formatting
] % Method A for two-column formatting

%%%%%%%%%%%%%%%  Main text   %%%%%%%%%%%%%%%
% \linenumbers

\section{Introduction}
% START-GRAPHENE PROPERIES %
Unique and remarkable quantum properties of monolayer graphene resembling a honeycomb structure, have sparked significant interest in both linear \cite{katsnelson} and nonlinear optics \cite{Agrawal}. % MONOLAYER GRAPHENE %
The Dirac cone symmetry is responsible for the ambipolar behaviour thus producing a surge of carrier concentration in suspended graphene \cite{ball1,ball2} and encapsulated graphene \cite{ball3}.% BALLISTIC ELECTRON TRANSPORT %
This is in contrast with the semiconductor where a particular quasi-particle is more significant than the other charge carriers. % CITE SC %
The charge carriers in graphene with zero rest mass resembles the relativistic entities and exhibiting an effective velocity comparable to the speed of light which are theoretically studied extensively \cite{mass1,mass2}. % HIGH MOBILITY %
These distinctive characteristics of charge carriers in pristine graphene are evident for its high conductivity \cite{rev_con1,rev_con2}. % hIGH COND.%

% TB %
In order to describe the electronic properties of different two dimensional materials, the tight-binding Hamiltonian  is a first simplified non-interacting model and commonly used on monolayer, bilayer, twisted graphene \cite{TBM_BG1,TBM_BG2,TBM_BG3} and TMDCs \cite{TBM_MOS1,TBM_MOS2}. This model reveals the behavior of $\pi$-electrons in the hexagonal lattice formed by carbon atoms. It considers the interaction between neighboring carbon atoms through electron hoping and overlapping of their atomic orbitals which successfully produces the key characteristics of graphene's electronic band structure including linear dispersion near the Dirac points.
% KUBO %
To understand the optical conductivity of graphene using the tight-binding Hamiltonian, one can employ the Kubo formula expression \cite{Kubo1,Kubo2,Kubo3} from the linear response theory \cite{mahan}. But current-current correlation function ($\Pi$) in the Kubo expression requires summations over all the states in higher dimensional systems making the calculations computationally expensive. Besides the Kubo approach,  the Semiconductor Bloch equation (SBE) \cite{mahan,ferry} is an alternative to calculate the optical conductivity of semiconductors and two dimensional crystals. Unlike Kubo theory, the SBE is derived from the time-dependent Schrödinger equation for a periodic systems which has great advantage in solving the out of equilibrium systems. 

In this paper, the linear optical response for monolayer graphene is investigated using the SBE approach with the many particle approximations at different doping levels. The Bloch states are described by a tight binding model and the expressions are limited to temperatures near zero Kelvin. The different scattering mechanisms are introduced with the aid of phenomenological relaxation parameters to the SBE and the conductivity response is obtained from the first order density matrix which comprises of both intra- and inter-band components.

Furthermore, the optical transmission analysis is correlated with experimental results to gain more insights about the Pauli blocking and conductivity response. Our formalism is in good agreement with the experimental data with slit deviations  due to the  approximations near Dirac points, lower optical field energy and the non-inclusion of thermal fluctuations.   We have shown that the utilization of approximations around Dirac points bridge the results obtained from perturbative calculations and the Kubo expression under an external electric field.

\section{Semi-classical Approach}\vspace{0.15 cm}

The SBE consists of three different types of Hamiltonian's which are  unperturbed, radiation-interaction and scattering to study the many-body physics. Unperturbed tight binding Hamiltonian deals with the nearest neighbour hopping with zero onsite energies.
    \begin{equation}
    \hat{H_0}=\sum\limits_{n}^{}\sum\limits_{k}\varepsilon_{nk}c_{nk}^{\dag}c_{nk}
    \end{equation}
The energy eigenvalues are denoted by $\varepsilon_{nk}$ with states n=\{1,2\} in the momentum space. The light-matter interacting Hamiltonian posses interaction between electric field of the light  and  charge distribution of matter. To construct the Hamiltonian, let's consider electric dipole approximation along with long wavelength limit. Under these approximations, there are two ways to write the light matter interaction Hamiltonian. One is using the velocity gauge given by \cite{sipe_gahramani},
 \begin{equation}
    H_{eA}=-\frac{e}{m}\left[\textbf{A}(t).\textbf{p}+e\textbf{A}^2/2\right]
    \end{equation}
and another one using the length gauge which has explicit dependence on external electric field given by,
  \begin{equation}  
    H_{eR}=-e\textbf{E}(t).\textbf{r}
    \end{equation}
 The main advantage of using the first expression with vector potential is preserving the crystal's translational symmetry, resulting the physical and geometrical response in k-space with the sum of independent contributions at each k point. The problem arises with formation of false divergences \cite{Peres17} obtained numerically at lower frequencies due to it's gauge sensitivity subjected to band truncation. Although, this problem can be diverted by the help of a proper complete basis set but solving for practical and numerical case is not possible, which leads to violation of gauge invariance due to bands truncation's. This band truncation has divergence issues when considering independent contributions from the omitted bands which produces less accurate results \cite{HTP2} over the length gauge methods. 
Sipe-Ghahramani \cite{sipe_gahramani} employed sum rules to deal with this divergence issues, they had to separate the interband and intraband contributions from the total Hamiltonian at the beginning which help them resolve the false errors. This was a key idea to deal with a quite similar problem in the length gauge analysis. Use of scalar potential has it's own perks due to presence of the gauge invariant electric field vector. We adopted this methodology due to advantages over velocity gauge. However, it does possess an additional scalar dependence that disturb the diagonal property of the perturbed Hamiltonian in k-space. Because of that, there will be some overlap or unnecessary coupling between k values thus producing 
 many divergences. Aversa $\&$ Sipe \cite{Aversa_Sipe} revisited this problem with their band isolation technique for linear response but it can also hold for subsequent higher orders. They observed that these spurious divergences actually have coefficients in which some of them can become zero. These sum rules seem to be violated during the use of approximations and that limits this methodology to length gauge.
     
The position operator in the length gauge approach is linked to the Berry connections when the bands are different from one another under external electrical filed. The electron wavefunctions can possibly exhibit a non local behaviour and thus spread over a number of sites. Close proximity of the Dirac cones (Fig. \ref{Energy_near_Dirac_points}) is no longer valid due to presence of external electric field radiation given by,
      \begin{equation}
      H_{eR}=-e\textbf{E}(t).\sum\limits_{n_1 n_2}\sum\limits_{k}c_{nk}^{\dag}\xi_{n_1n_2k}c_{nk}-ie\textbf{E}(t).\sum\limits_{n}^{}\sum\limits_{k}\nabla_{k}c_{nk}^{\dag}c_{nk}
      \end{equation}
Second term in the above equation, possesses local band characteristics. The effect of electron scattering due to phonon, impurities and electrons is dealt by the scattering Hamiltonian. In the equation of motion, the rate of change due to scattering can be made proportional to the density of states with the inclusion of a  phenomenological relaxation parameters for any order.
Combining all the three sub Hamiltonians from above leads to the construction of total Hamiltonian.
 
\begin{equation}
\hat{H}=\hat{H_0}+\hat{H}_{eR}+\hat{H}_{sca}  \label{total}
\end{equation} 
    
% Figure environment removed 

\subsection{Collisional Multiband Boltzmann Equation}
\vspace{0.15 cm}
In this paper, we have adopted Sipe-Aversa's SBE \cite{Aversa_Sipe} with length gauge perturbation theory to  get the equation of motion Eq.(\ref{SBE}) and solved for the response function of monolayer graphene. Cheng, Sipe and Vermuelen carried out an extensive study on doped monolayer graphene to obtain linear and nonlinear optical conductivity with scattering \cite{CSV15} and without scattering \cite{CSV141} using approximations around Dirac Points. Using similar assumptions, they have calculated DC Current induced Second harmonic generation \cite{CSV142} at zero temperature limit. A quite similar research has been conducted by Hipolito,  Taghizadeh, Alireza and Pedersen \cite{HTP1} to find the linear conductivity, third harmonic generaion and optical Kerr effect for both monolayer and bilayer graphene. Identification of new kind of third order divergences during cross phase modulation and degenerate four wave mixing were carried out by Cheng, Sipe, Wu and Guo \cite{cheng_intraband} for centrosymmetric two dimensional materials. Naib and Sipe \cite{Naib} detected intraband and interband current with applying electric field through numerical simulation on undoped monolayer graphene at low temperatures. McGouran et al \cite{mcgouran} obtained dipole matrix elements through length gauge and found the linear and nonlinear Terahertz responses for undoped suspended bilayer graphene through numerical simulations at different low temperatures with scattering time of 80 femtoseconds. Study of carrier mobility, carrier density, band structure, linear and nonlinear TeraHertz response of nitrogen doped graphene were carried out by Anvari et al \cite{nitrogen}. Cheng and Guo \cite{magnetic} studied the nonlinear magneto-optic effects in doped and gapped graphene. They examined linear response, THG, Kerr effects, two photon absorption and four wave mixing under external magnetic field (B=0.05 T). In general, most of the calculations are always focused on interband scattering contribution while neglecting the intraband scattering.Theoretically, the length gauge approach for  equation of motion with the density matrix is used for an innumerable times to investigate various optical parameters. Using modified SBE, we takes intraband information directly into account with the aid of phenomenological relaxation parameters. To obtain expression for many body interactions (Eq.(\ref{CMBE})), we followed semi-classical approach using the quantum Liouville equation,
\begin{equation}
    i\hbar\frac{\partial \hat{\rho}(t)}{\partial t}=[\hat{H},\hat{\rho}(t)]
\label{liou}
\end{equation}
which consists of the total Hamiltonian Eq.(\ref{total}) and the time dependent density operator,
\begin{equation}
\hat{\rho}(t)=\sum\limits_{n_1 n_2}^{}\rho_{n_1 n_2}(t) |n_1\rangle \langle n_2|
\label{density}
\end{equation}
Using Eq. (\ref{total}),(\ref{liou}),(\ref{density}) the equation of motion in terms of density matrix is obtained,
\begin{multline}
i\hbar\frac{\partial {\rho_{n_1 n_2}}(t)}{\partial t}|n_1\rangle \langle n_2|= H_0 |n_1\rangle\rho_{n_1 n_2}(t) \langle n_2|-\rho_{n_1 n_2}(t) |n_1\rangle\langle n_2|H_0 \\
-e\textbf{E}.[\textbf{r},\rho_{n_1 n_2}(t)] |n_1\rangle\langle n_2|-i\Gamma_{n_1n_2}\rho_{n_1n_2}(t)|n_1\rangle\langle n_2|
\label{SBE}
\end{multline}
 where, $\Gamma$ is the phenomenological constant which contains the scattering parameters. In the presence of the scattering,  the time evolution of the density operator undergo various physical effects including damping or decoherence. We have used a Blount equation \cite{blount} to get a commutation relationship for the spatial coordinate (\textbf{r}) with a physical operator($\Lambda$),
\begin{equation}
\int{dr\Psi_{n_1k_1}^{*}(r)[r,\Lambda]\Psi_{n_2k_2}(r)}=[\xi_{k_1k_2},\Lambda_{k_1k_2}]_{n_1,n_2}+i\nabla_k\Lambda_{n_1n_2k_1k_2}
\label{Blount}
\end{equation}

If this Blount formula is extended for each k values, we get the required collisional multiband Boltzmann equation (CMBE),
\begin{multline}
i\hbar\frac{\partial {\rho_{n_1 n_2k}}(t)}{\partial t}=(\varepsilon_{n_1k}-\varepsilon_{n_2k})\rho_{n_1 n_2k}(t)-i\Gamma_{n_1n_2k}\rho_{n_1n_2k}(t)+\\      -e\textbf{E}(t).\sum\limits_{n}\left(\xi_{n_{1}nk}\rho_{nn_2k}- \rho_{n_1nk}\xi_{nn_2k}\right)-ie\textbf{E}(t).\frac{\partial\rho_{n_1n_2k}}{\partial k}
\label{CMBE}
\end{multline}

\subsection{Conductivity Expression}
\vspace{0.15 cm}
Composition of the j$^{th}$ order CMBE is possible, if the density matrix is considered to be sum of all orders
    \begin{multline}
    \hbar \frac{\partial \rho_{n_1n_2k}^{\!(j)}(t)}{\partial t}=-i\left[\varepsilon_k,\rho_k^{\!(j)}{\!(t)}\right]_{n_1n_2}-\Gamma_{n_1n_2k}^{\!(j)}\rho_{n_1n_2k}^{\!(j)}(t)+\\ ie\textbf{E}(t).\left[\xi_k,\rho_k^{\!(j-1)}\right]_{n_1n_2}-eE(t).\nabla_k\rho_k^{\!(j-1)}
\label{jSBE}
    \end{multline}
The solution of the CMBE equation is in the following form,
\begin{equation}
    \rho_{n_1n_2k}^{\!(j)}=\int..\int\left[-e\frac{d\omega}{2\pi}E_\omega^ae^{-i\omega t}\right]^{j}\mathcal{S}_{n_1n_2k}^{\!(j)}(\omega_1,\omega_2,..,\omega_j) 
\label{solution}
\end{equation}
The j$^{th}$ frequency term can be procured by simple Fourier transform which contains the inter/intra-band Berry connection elements. The perturbative expansion of the current density can be used to equate the areal current density,
    \begin{multline}
       e\sum\limits_{n_1n_2}\int\frac{d\textbf{k}}{4\pi^2}v^a_{n_1n_2k}\rho_{n_1n_2k}^{\!(j)}(t)\\
       =\int..\int\left[\frac{d\omega}{2\pi}E_\omega^ae^{-i\omega t}\right]^{j}\sigma^{\!(j);abcd..}(\omega_1,\omega_2,.., \omega_j)
       \label{current_density}
    \end{multline}
    After substitution of Eq.(\ref{solution}) in Eq.(\ref{current_density}), we get the the general formula of conductivity,
    \begin{equation}
    \sigma^{\!(j);abcd..}=-e^{j+1}\sum_{n_1n_2}\int\frac{d\textbf{k}}{4\pi^2}v^a_{n_1n_2k}\mathcal{S}_{n_1n_2k}^{\!(j);bcd..}(\omega_1,\omega_2,.., \omega_j)
    \label{formula_conductivity}
    \end{equation}
In connection to this, direct gap approximation is needed to be taken into the account when the Berry connection is dealing with interband transition. It is well assumed that the momentum states in both bands (CB and VB) has to be aligned ($\Delta k=0$), thus restricting the oblique transitions. Not only, for Berry connection this approximation is true to any light matter coupling operator whose matrix element is defined as
\begin{equation}\left\langle ck|\hat{h^{'}}|vk\right\rangle=\int{dr\Psi_{ck}^{*}(r)\hat{h^{'}}\Psi_{vk}(r)} \end{equation}
As long as the external field or the applied perturbation is of the same order as the measured response,we are in the linear regime.  
From the  CMBE, the first order equation is deduced as follows,
\begin{multline}
\hbar \frac{\partial \rho_{n_1n_2k}^{\!(1)}(t)}{\partial t}=-i\left[\varepsilon_k,\rho_k^{\!(1)}{\!(t)}\right]_{n_1n_2}-\Gamma_{n_1n_2}^{\!(1)}\rho_{n_1n_2k}^{\!(1)}(t)\\+ie\textbf{E}(t).\left[\xi_k,\rho_k^{\!(0)}\right]_{n_1n_2}-eE(t).\nabla_k\rho_k^{\!(0)}
\label{1st_CMBE}
\end{multline}
The solution in Eq.(\ref{solution}) is found to be of usual form whose time derivative has only one time dependent term,
\begin{equation}
\frac{\partial \rho_{n_1n_2k}^{\!(1)}(t)}{\partial t}=-i\omega_1\rho_{n_1n_2k}^{\!(1)}(t)
\label{derivative}
\end{equation}
If Eq.(\ref{derivative}) is substituted in the first order CMBE (Eq.(\ref{1st_CMBE})), the expression becomes,
\begin{multline}
-i\hbar\omega_1\rho_{n_1n_2k}^{\!(1)}(t) =
-i\hbar\omega_{n_1n_2k}\rho_{n_1n_2k}^{\!(1)}(t)-\\ \Gamma_{n_1n_2}^{\!(1)}\rho_{n_1n_2k}^{\!(1)}(t)+ie\textbf{E}(t).\left[\xi_k,\rho_k^{\!(0)}\right]_{n_1n_2}-eE(t)\frac{\partial}{\partial k}\rho_{n_1n_2k}^{\!(0)}
\end{multline}
where, the two band frequency is actually the difference between any two bands ($\omega_{n_1n_2}=\omega_{n_1}-\omega_{n_2}$) and the phenomenological constant is actually dependent on the signs of $n_1$ and $n_2$.
Eventually, taking the Fourier transform of electric field one can find the final form just by simple comparison of the density matrix element as prescribed above.
\begin{equation}
\rho_{n_1n_2k}^{\!(1)}(t)=\frac{(-e)\bigintsss{\frac{d\omega_1}{2\pi}E_{\omega_1}^ce^{-i\omega_3t}}\left(\left[\xi^c_k,\rho_k^{\!(0)}\right]_{n_1n_2}+\frac{\partial}{\partial k_c}\rho_{n_1n_2k}^{\!(0)}\right)}{\hbar\omega_1-\hbar\omega_{n_1n_2k}+i\Gamma_{n_1n_2}^{\!(1)}}
\end{equation}
The first order amplitude term \cite{cheng_book} is given by,
\begin{equation}
\mathcal{S}_{n_1n_2k}^{\!(1);c}(\omega_1)=\frac{\left[\xi^c_k,\rho_k^{\!(0)}\right]_{n_1n_2}+\frac{\partial}{\partial k_c}\rho_{n_1n_2k}^{\!(0)}}{\hbar\omega_1-\hbar\omega_{n_1n_2k}+i\Gamma_{n_1n_2}^{\!(1)}}
\label{amplitude}
\end{equation}
From the above expression of linear amplitude one can separate out the interband part (first term) and intraband component (second term ).  Then Eq.(\ref{amplitude}) is substituted  in Eq.(\ref{formula_conductivity}) to get the first order conductivity expression which is eventually a crucial step in matching the Kubo results (Appendix B). 
%The first term in Eq.(\ref{amplitude})  being concerned with the Berry connection, will give a non zero value if and only if the signs of $n_1$ and $n_2$ are different. The other term is simply associated with the diagonal elements of the usual 2x2 density matrix.

\section{Results}

Graphene exhibits unique linear optical responses, particularly when it is subjected to external stimuli.
These fascinating properties are due to two dimensional confinement and behaviour of charge carriers near Dirac points. In order to explore some of these unusual cases, the equations in Eq.(B.4) and Eq.(B.9), are used for calculating the real and imaginary part of the optical conductivity.
% Figure environment removed
The real component of the linear response term without scattering or thermal effects will consist only the interband contribution as a simple step function and shows optical shielding behaviour which is consistent with reported literature \cite{CSV141}. Now, with the inclusion of finite scattering ($\Gamma_e$ for interband as in Eq.(B7) and $\Gamma_i$ for intraband as in Eq.(B2)), the terms are expressed separately with the rejection of logarithmic term. The scattering values are taken from the experimental results, $\Gamma_e=4.14$ meV, $\Gamma_i=0.414$ meV \cite{Malard} and $\Gamma_e=0.5$ meV, $\Gamma_i=65$ meV \cite{gu}.
%and real, imaginary part of the conductivities are plotted in the Fig. \ref{conductivity}. 
A vital observation from Fig.\ref{Drude} is showcased with respect to Drude response for different values of chemical potential and scattering rates. The inter-band scattering is dominating due to the large availability of free charge carries which is responsible for the strong Drude response as shown in the Fig.\ref{Drude}(a). Whereas in the Fig.\ref{Drude}(b), it clearly shows that the intraband scattering emerges as a new phenomena which reduces the Drude response drastically and increases its bandwidth. We clearly observed that Drude behavior is tunable with the interplay between interband and intraband scattering. Further we examine the real conductivity response at different frequency regimes where Pauli Blocking (PB) window starts to appear for different scattering rates under the influence of doping. Fig.\ref{conductivity}(a) is demonstrated with respect to PB effect which is found to be more effective or sharper with the reduction of interband scattering rate $(\Gamma_e=10 \Gamma_i)$ while the slope of the PB curve decreases leading to radiation leakage with higher scattering rates $(\Gamma_e=130 \Gamma_i)$ and similar kind of behavior is observed by Hipolito et al \cite{HTP1}. Conductivity ratio value reaches unity in both cases at different frequencies depending upon scattering parameters.

When we reverse the scattering dominance to intraband type, we observed that the Drude behavior is suppressed and there is emergence of  new conductivity peak at the beginning of PB which is illustrated in Fig.\ref{conductivity}(c). Surprisingly, the conductivity emerges beyond the universal conductivity due to dominant intraband scattering with these set of values $\Gamma_e = 0.5$ meV and $\Gamma_i = 65$ meV \cite{gu}. This feature is consistent for different doping values but the peak amplitude decreases exponentially (red line in Fig.\ref{enhance}) with increasing doping concentration. A noteworthy aspect is that the aforementioned unique feature appears from a certain doping threshold value and below this value free carrier Drude response dominates. The higher intraband scattering dominance in our case is closely matching  with the experimental results obtained by Basov et al. at specific external bias condition \cite{Henriksen}. These experiments were conducted at 28 V gate bias voltage (Fermi level at 0.18 eV) at 45 K, in which the PB curve is bending in the conductivity plot from 0.223 eV to 0.447 eV. Whereas in our case, calculations have produced a sharp change in response at 0.367 eV (line \textbf{A}) at zero Kelvin which is lying in region where the slope reaches the maximum as depicted in Fig.\ref{conductivity}(c). The Auger recombination (AR) and impact excitation (IE) are influenced by the PB which eventually mediates the excitonic effects \cite{exciton_1, exciton_2}. These electrically neutral quasi-particles participate in band renormalization and compete with the PB mechanism which may produce such unique optical responses during negligible Coulomb screening.  

The imaginary part of the conductivity for different scattering rates and doping levels are plotted in Fig.\ref{conductivity}(b,d). The absorption peaks are logarithmic divergences around PB in the conductivity curves.
It is evident that the peak height increases with doping and the peak amplitude gets quenching for a high scattering rate. For direct band gap two dimensional materials, the energy levels of carriers are shifted due to the interactions between electrons and impurities, it will enhance the absorption rates at higher doping levels \cite{Lahourpour_2019}. With the variation of chemical potential one can effectively shift the absorption peak to a desired frequency.  From Fig.\ref{conductivity}(d), the peak at 0.2 eV correspond to Mid-Infrared Range is smoothly shifted to  Near-Infrared Range with slight increment of chemical potential. For a higher interband scattering $(\Gamma_e=130\Gamma_i)$ with respect to intraband scattering the broadening of absorption peak occurs.It should be noted that the exact behaviour can be obtained through Random Phase Approximation (RPA) calculations \cite{Koppens_2011} with the aid of a finite temperature extension at 300 K and relaxation rate of $10^{-13}$ secs. The smearing for the real term and the broadening of the imaginary part is mainly caused by the consideration of thermal effects but those are inherently present in our calculations.
% Figure environment removed

% Figure environment removed


\section{Conclusion}
\vspace{0.2cm}
We have conducted a systematic study on linear optical properties of monolayer graphene under low doping conditions to determine the dominant mechanism and interplay between interband $\&$ intraband scattering. 
%Constructed collisional multiband Boltzmann equation with the use of length gauge for studying many body effects.
A collisional multiband Boltzmann equation was constructed using the length gauge for understanding many-body effects. 
%Drude behavior is tunable with phenomenological scattering parameters $(\Gamma_e, \Gamma_i)$ and high amplitude quenching with broadening of bandwidth occurs for interband scattering event.
The Drude behavior can be tuned through phenomenological scattering parameters and specifically for intraband scattering, a high-amplitude quenching along with broadening of the bandwidth. During the dominance of intraband scattering  under low doping regime, we have observed the emergence of new conductivity peak which is higher than the universal conductivity for after Pauli blocking window. The leap in response by excitonic effects in monolayer graphene is primarily triggered by dominant intraband scattering rather than interband scattering. It is very crucial that the material parameters of graphene, such as carrier density, temperature, and impurity concentration, will dictate the importance of different scattering processes. 
Our numerical study shed light on the monolayer graphene to explore the novel optical response and control on pauli blocking region with doping concentrations for strong optical absorption.  





\section*{Acknowledgment}\vspace{0.15 cm}
We greatfully acknowledge support from DST Science and Engineering Research Board (SERB), India under grant No. SERB:CRG/2022/008749 and  IIT Jodhpur Seed grant No.I/SEED/BMK/20230017.  

\section*{Appendix}
%\appendix
\appendix
\renewcommand{\theequation}{A.\arabic{equation}}
\subsection*{Appendix A: Derivation of Results around Dirac Point}

%\subsection{} 

\setcounter{equation}{0}
\label{app:section1}
% Approximations around Dirac Point: APPENDIX-A
%Derive linear energy term
The tight binding wavefunction consists of set of all 2$p_z$ orbitals for a particular state with $n=2$ atomic orbitals in the unit cell, given by
\begin{equation} 
\phi_A(r)=\frac{1}{\sqrt{N}}\sum\limits_{R}e^{ik.R}\varphi_A(r-R)
\label{A1}
\end{equation}
The normalisation constant (N) is the total number unit cells present per unit area and these Wannier orbitals $\varphi(r-R)$ are well localised around respective atomic sites (R) of Bravais lattice. To ensure periodicity another parameter ($\tau$) is included with the position which relates site of one sublattice to that of second sublattice. If it is absent for one site, it must be non zero for the other one.
The wavefunction corresponding to the electrons in real crystal can be written in terms of tight binding wavefunctions for two inequivalent sites A and B  denoted by $l$ below \cite{tightbinding}. The band indices are hidden to make the situation simpler for now-
\begin{equation}
    \Psi_k(r)=a_{k}\phi_A(r)+b_{k}\phi_B(r)
             =\sum\limits_{l=1}^2C_{lk}\phi_l(r)           
\end{equation}
Using Eq.(\ref{A1}), the wavefunction is written as,
\begin{equation} \Psi_k(r)=\frac{1}{\sqrt{N}}e^{ik.r}\sum\limits_{R}e^{ik.(R-r)}\sum\limits_{l=1}^2\varphi_i(r-R)C_{lk}=\frac{1}{\sqrt{N}}e^{ik.r}u_k(r)
\label{periodic_term}
\end{equation}\\
This newly formed state $\Psi_k(r)$ with the periodic term $u_k(r)$ is the Bloch function used as eigenvectors for the unperturbed Hamiltonian.The main approximation of the tight binding model is to present the respective Bloch functions with the aid of well localised atomic orbitals which is successively worked out above. The Fourier transform of creation and annihilation operators is used for k-space transformation as shown below, 
\begin{multline}
\hat{H_0}=\sum\limits_{i}\varepsilon_Aa_i^{\dag}a_i+\sum\limits_{j}\varepsilon_Bb_j^{\dag}b_j\\
+\sum\limits_{<i,j>}t_{ij}a_i^{\dag}b_j+\sum\limits_{<i,j>}t_{ij}b_j^{\dag}a_i
     =\sum\limits_{k}c_k^{\dag}H_{ij}(k)c_k
     \label{A.4}
     \end{multline}
 The transfer integral matrix for monolayer graphene from Eq.(\ref{A.4}) with zero onsite energy is,  
    \begin{equation}{H_{ij}}=
     \begin{pmatrix}0 & \gamma_0s(k) \\\gamma_0s^{*}(k) & 0 \end{pmatrix}
     \label{A.5}
     \end{equation}
where, $c_{nk}^{\dag}$ and $c_{nk} $ are electron operators which are linear combinations of $a_k$ and $b_k$.
In the above equation, s(k) is the structure factor dependent upon positions in real space .
\begin{equation}
  s(k) =\sum\limits_me^{ik.\delta_m} = e^{ik_y a}+2e^{-i\frac{k_y a}{2}} \cos{\frac{\sqrt{3}}{2}k_xa}
\end{equation}
Corresponding measure of distance for particles in real space is denoted by $\delta_m $ and -t is the hopping matrix equivalent to $\gamma_0$ respectively. If the Dirac point K(K') is given by $(\pm4\pi/{3\sqrt{3}a},0)$, the structure factor in the neighborhood of K' can be written as, 
\begin{equation}
  s(K'+k) =e^{-ik_ya}+2e^{-i\frac{k_y a}{2}} \cos\left(\frac{\sqrt{3}}{2}k_xa-\frac{2\pi}{3}\right)
\end{equation}
To the lowest order in k, we achieve the final form,
\begin{equation}
  s(K'+k) =\frac{3a}{2}(k_x+ik_y)
\end{equation}
After performing diagonalization on Eq.(\ref{A.5}), we obtain the following eigenvalues,

\begin{equation}
    \varepsilon_{nK'+k}=\gamma_0|s(K'+k)|=\frac{3a\gamma_0}{2}|k|
\end{equation}
%From Figure A1, if we consider point $A_1$ as the origin, the distances (red ones) are $\delta_1$=(-a/$2\sqrt{3}$,-a/2), $\delta_2$=(a/$\sqrt{3}$,0), $\delta_3$=(-a/$2\sqrt{3}$,a/2). 
This is the case when there is negligence of overlap between orbitals on different sites and based on the NN approximation only the electron hopping is allowed for one of the three neighbouring lattice sites. The energy is stated in a compact manner with the substitution of Fermi velocity $v_F=(3a\gamma_0/2\hbar) $ finally leading to a linear dispersion relation near the Dirac points,
\begin{equation}
    \varepsilon_{nK+k}=\pm \hbar v_F k
\end{equation}
%\section{Velocity Components}
As per requirement for any inter or intra-band parts we must start to mention the band indices n which is crucial for bifurcating the velocity components.  Due to the orthonormal nature of the wavefunction, it follows these following conditions,\\\\
Normalisation condition: 
\begin{equation}
\int{dr\Psi_{n_1k_1}^{*}(r)\Psi_{n_2k_2}(r)}=\delta_{n_1n_2}\delta(k_1-k_2)
\label{A11}
\end{equation}
Matrix element formula:
\begin{equation}
    \int{dr\Psi_{n_1k_1}^{*}(r)A\Psi_{n_2k_2}(r)}=A_{n_1n_2k_2}\delta(k_1-k_2)
    \label{A12}
    \end{equation}
The Hamiltonian matrix component is given by,
\begin{align*}
    H_{n_1n_2k_1k_2}&=\int{dr\Psi_{n_1k_1}^{*}(r)H\Psi_{n_2k_2}(r)}\\
    &=\varepsilon_{n_2k_2}\int{dr\Psi_{n_1k_1}^{*}(r)\Psi_{n_2k_2}(r)}
  \end{align*} 
  Applying Eq.(\ref{A11}) in the above equation,
   \begin{equation}
    H_{n_1n_2k_1k_2}=\varepsilon_{n_1k_1}\delta_{n_1n_2}\delta(k_1-k_2) 
    \end{equation}

Now, let's simplify for the first half of the Blount expression.
\begin{equation}
    [\xi_{k_2},H_{k_1k_2}]_{n_1n_2}=\left(\varepsilon_{n_2k_2}-\varepsilon_{n_1k_1}\right)\int{dr\Psi_{n_1k_1}^{*}(r)\xi_{k_2}\Psi_{n_2k_2}(r)}
\end{equation}
Using Eq.(\ref{A12}), the commutation relation becomes,
\begin{equation}
[\xi_{k_2},H_{k_1k_2}]_{n_1n_2}=\left(\varepsilon_{n_2k_1}-\varepsilon_{n_1k_1}\right)\xi_{n_1n_2k_1}\delta(k_1-k_2)
\label{A14}
\end{equation}
Thus we have all the equations which are necessary for the construction of the velocity matrix element. We will start with the Heisenberg equation of motion followed by the use of Eq.(\ref{Blount}), Eq.(\ref{A12}) and Eq.(\ref{A14}) respectively,
\begin{align*}
\int{dr\Psi_{n_1k_1}^{*}(r)v\Psi_{n_2k_2}(r)}&=\int{dr\Psi_{n_1k_1}^{*}(r)\frac{1}{i\hbar}[r,H]\Psi_{n_2k_2}(r)}\\
&=\frac{1}{i\hbar}\left([\xi_{k_2},H_{k_2}]_{n_1,n_2}+i\nabla_kH_{n_1n_2k_2}\right)
%&=\frac{-i}{\hbar}\left(\left(\varepsilon_{n_2k_1}-\varepsilon_{n_1k_1}\right)\xi_{n_1n_2k_1}+i\nabla_k\varepsilon_{n_1k_1}\delta_{n_1n_2}\right)\delta(k_1-k_2)
\end{align*}
\begin{equation}
\int{dr\Psi_{n_1k_1}^{*}(r)v\Psi_{n_2k_2}(r)}=v_{n_1n_2k_1}\delta(k_1-k_2)
\end{equation}
The intraband term $(n_1=n_2)$ is then given by $\frac{1}{\hbar}\nabla_k\varepsilon_{n_1k_1}\delta_{n_1n_2}$ and the interband term $(n_1\neq n_2)$ is $i\omega_{n_1n_2k_1}\xi_{n_1n_2k_1}$ composedly giving \cite {mikhailov},
\begin{equation}
v_{n_1n_2k_1}=\frac{1}{\hbar}\nabla_k\varepsilon_{n_1k_1}\delta_{n_1n_2}+i\omega_{n_1n_2k_1}\xi_{n_1n_2k_1}
\end{equation}
 
%\section{Berry Connection}
After the evaluation of velocity matrix elements, the next exercise is solely focused to express Berry connections in a second quantised notation which can be further simplified to evaluate the approximation done near K points for the interband part.
Blount's conjecture of an irrational meaning of position matrix is obeyed here. Specifically, in the context of condensed matter physics, the Berry's connection is related to the area of the unit cell through the expression for the Berry phase. If the subscript "uc" in area $A_{uc}$ signifies the unit cell of the crystal lattice, then the connection formula is given by \cite{berry},
\begin{equation}\xi_{n_1n_2k}=i\int{\frac{dr}{A_{uc}}u_{n_1k}^{*}(r)\nabla_ku_{n_2k}(r)}\end{equation}
As defined earlier in Eq.(\ref{periodic_term}), the periodic part of Bloch function is found to be 
\begin{equation}u_{nk}(r)=\sum\limits_{R}e^{ik.(R-r)}\left[\varphi(r-R-\tau_A)a_{k}^n+\varphi(r-R-\tau_B)b_{k}^n\right]
\label{Periodic_second}
\end{equation}
Some general rules followed by the Wannier orbitals which will be used in the next steps are listed below,\\\\
Normalisation condition:
\begin{equation}
\int{\varphi^{*}(r-R-\tau_A)\varphi(r-R-\tau_A)dr}=1
\label{A.19}
\end{equation}
Orthogonal condition:
\begin{equation}\int{\varphi^{*}(r-R-\tau_A)\varphi(r-R-\tau_B)dr}=0\end{equation}
Meaninglessness of space matrix element:
\begin{equation}\int{\varphi^{*}(r-R-\tau_A)(r-R-\tau_A)\varphi(r-R-\tau_A)dr}=0
\label{A.21}
\end{equation}
Let's begin with the gradient of $u_{n_2k}(r)$
\begin{multline} 
\nabla_ku_{n_2k}(r)=\sum\limits_{R}i(R-r)e^{ik.(R-r)}\sum\limits_{R}e^{ik.(R-r)}[\varphi(r-R\\
-\tau_A)a_{k}^{n_2}+\varphi(r-R-\tau_B)b_{k}^{n_2}]+\sum\limits_{R}e^{ik.(R-r)}\varphi(r-R\\
-\tau_A)\nabla_ka_{k}^{n_2}
+\sum\limits_{R}e^{ik.(R-r)}\varphi(r-R-\tau_B)\nabla_kb_{k}^{n_2}
\label{A.22}
\end{multline}
The final expression will be a 2x2 matrix with the off diagonal elements getting null values due to the orthogonal relation.
So only, the diagonal elements [AA,BB] has to be disentangled. Placing Eq.(\ref{Periodic_second}) and (\ref{A.22}) in the connection formula, there should be two terms remaining for each of the diagonal elements, 
\begin{multline}
i\int{\frac{dr}{A_{uc}}u_{n_1k}^{*}(r)\nabla_ku_{n_2k}(r)}|_{AA_1}=
\frac{1}{A_{uc}}a_{k}^{n_1\dag}\sum\limits_{R'R}e^{-ik.(R'-r)}e^{ik.(R-r)}\\
\int\varphi^{*}(r-R'-\tau_A)
(r-R)\varphi(r-R-\tau_A)dr a_k^{n_2}
\end{multline}
From Eq.(\ref{A.21}), the above equation can be written as,
\begin{multline}
    i\int{\frac{dr}{A_{uc}}u_{n_1k}^{*}(r)\nabla_ku_{n_2k}(r)}|_{AA_1}
    =\frac{1}{A_{uc}}a_{k}^{n_1\dag}\\\sum\limits_{R}
    \int{\varphi^{*}(r-R-\tau_A)\tau_A\varphi(r-R-\tau_A)dr}a_k^{n_2}
    \end{multline}
Using Eq.(\ref{A.19}), it can be further simplified to,
    \begin{equation}
   i\int{\frac{dr}{A_{uc}}u_{n_1k}^{*}(r)\nabla_ku_{n_2k}(r)}|_{AA_1}
   =a_{k}^{n_1\dag}\tau_Aa_k^{n_2}
   \label{A.25}
    \end{equation}
\begin{multline}
    i\int{\frac{dr}{A_{uc}}u_{n_1k}^{*}(r)\nabla_ku_{n_2k}(r)}|_{AA_2}=\frac{i}{A_{uc}}a_{k}^{n_1\dag}\\
    \sum\limits_{R}\int{\varphi^{*}(r-R-\tau_A)\varphi(r-R-\tau_A)dr}\nabla_ka_k^{n_2}\end{multline}
    Again, applying Eq.(\ref{A.19}) to the second element,
    \begin{equation}
    i\int{\frac{dr}{A_{uc}}u_{n_1k}^{*}(r)\nabla_ku_{n_2k}(r)}|_{AA_2}=a_{k}^{n_1\dag}i\nabla_ka_k^{n_2}
    \end{equation}
If the exact exercise is done for the B sub-lattice we get,
\begin{equation} i\int{\frac{dr}{A_{uc}}u_{n_1k}^{*}(r)\nabla_ku_{n_2k}(r)}|_{BB_1}=a_{k}^{n_1\dag}\tau_Bb_k^{n_2}b_k^{n_2}\end{equation}
\begin{equation} i\int{\frac{dr}{A_{uc}}u_{n_1k}^{*}(r)\nabla_ku_{n_2k}(r)}|_{BB_2}=b_{k}^{n_1\dag}i\nabla_kb_k^{n_2}
\label{A.28}
\end{equation}
Coupling Eq.(\ref{A.25}-\ref{A.28}) in Eq.(\ref{Periodic_second}), we can easily obtain,
\begin{equation}  \xi_{n_1n_2k}=a_{k}^{n_1\dag}(\tau_A+i\nabla_k)a_k^{n_2}+b_{k}^{n_1\dag}(\tau_B+i\nabla_k)b_k^{n_2}
\label{final_berry}
\end{equation}



The required matrix element included in the dipole interaction Hamiltonian is given in second quantised format for each sublattice. It states the connection between two neighbouring wavefunctions which is applicable for both intraband and interband cases (Fig.\ref{berry_plot})
\begin{equation} 
\xi_{n_1n_2k}=\begin{pmatrix}a_{k}^{n_1\dag} & b_{k}^{n_1\dag}\end{pmatrix}\begin{pmatrix}\tau_A+i\nabla_k & 0 \\0 & \tau_B+i\nabla_k \end{pmatrix}\begin{pmatrix}a_k^{n_2} \\b_k^{n_2} \end{pmatrix}
\end{equation}
where, the eigenvectors from Eq.(\ref{A.4}) are given by,
\begin{equation}\begin{pmatrix}a_k^{n_2} \\b_k^{n_2} \end{pmatrix}=\frac{1}{\sqrt2}\begin{pmatrix}n\frac{s(k)}{|s(k)|} \\1 \end{pmatrix}\end{equation}

% Figure environment removed
\renewcommand{\theequation}{B.\arabic{equation}}
\subsection*{Appendix B: Linear Conductivity Expression} 
\setcounter{equation}{0}
\label{app:section2}
The intraband term is derived below in the neighborhood of the zero Kelvin temperature scale with the energy terms (from Appendix A) written as, $\varepsilon_{1k}=+\hbar v_Fk=\varepsilon$ and $ \varepsilon_{2k}=-\hbar v_Fk=-\varepsilon$. Similarly, the velocity terms are given by, $
 v_{11k}^x=+v_{F}\frac{k_x}{k}$ and $ v_{22k}^x=-v_{F}\frac{k_x}{k}$.
 From Eq.(\ref{formula_conductivity}) the intraband linear conductivity expression is written as,
 \begin{equation}
\sigma_{intra}^{\!(1)xx}(\omega) =-e^2\int{\frac{d\textbf{k}}{4\pi^2}v_{11k}^x\mathcal{S}_{11k}^{ \!(1);x}(\omega)}-
e^2\int{\frac{d\textbf{k}}{4\pi^2}v_{22k}^x\mathcal{S}_{22k}^{ \!(1);x}(\omega)}     
\end{equation}
The intraband frequency coefficient values are now placed in the above expression along with the approximations taken around the Dirac point,
 \begin{multline}
\begin{split}
\sigma_{intra}^{\!(1)xx}(\omega) = -e^2\int{\frac{d\textbf{k}}{4\pi^2}} \left[\frac{iv_{11k}^x}{(\hbar\omega+i\Gamma_i^{\!(1)})}\left(\frac{\partial f(\varepsilon)}{\partial \varepsilon}\frac{\partial \varepsilon}{\partial k_x}\right) \right. \\
\left. +\frac{iv_{22k}^x}{(\hbar\omega+i\Gamma_i^{\!(1)})}\left(\frac{\partial f(-\varepsilon)}{\partial \varepsilon}\frac{\partial \varepsilon}{\partial k_x}\right) \right]
\end{split}
\label{B2}
\end{multline}
 where, $f(+\varepsilon)$ or $f_{1k}$ is the Fermi-Dirac Distribution function. \\
If the intraband velocity  matrix elements are substituted in the numerator, the expression can be rewritten as,
\begin{multline}
\sigma_{intra}^{\!(1)xx}(\omega) = ie^2\int{\frac{d\textbf{k}}{4\pi^2} \left(\frac{k_x}{k}\right)^2 \frac{v_F}{k}} \\
 \frac{hv_Fk}{(\hbar\omega+i\Gamma_i^{\!(1)})} \left(\frac{\partial f(-\varepsilon)}{\partial\varepsilon}-\frac{\partial f(\varepsilon)}{\partial\varepsilon}\right)
\end{multline}
The transformation from momentum space to energy space is carried out to produce one half of the result obtained by Kubo expression,
\begin{equation}
\sigma_{intra}^{\!(1)xx}(\omega) =\frac{ie^2}{\pi\hbar(\hbar\omega+i\Gamma_i^{\!(1)})}\int\limits_0^{\infty}\varepsilon\left(\frac{\partial f(-\varepsilon)}{\partial\varepsilon}-\frac{\partial f(\varepsilon)}{\partial\varepsilon}\right)d\varepsilon
\label{Kubo_intra}
\end{equation}
The above conversion can also be carried out using an identity which is only applicable for two-dimensional isotropic materials given by, 
\begin{equation}
    \int{\frac{d\textbf{k}}{4\pi^2} \left(\frac{k_x}{k}\right)^2 \frac{v_F}{k}}=  \int{\frac{d\textbf{k}}{4\pi^2} \left(\frac{k_y}{k}\right)^2 \frac{v_F}{k}}=\frac{1}{\pi \hbar}\int\limits_0^{\infty}{d\varepsilon}
\end{equation}
The cold semiconductor approximation is put forward to achieve close to zero temperature results.  An explicit transformation of chemical doping can be put forward when temperature is finite. 

Similarly, for the interband term, the required velocity approximations are listed below,
$v_{12k}^x=iv_{F}\frac{k_y}{k}$ and $v_{12k}=-iv_{F}\frac{k_y}{k}$
along the Berry connection elements $\xi_{12k}^x=\frac{k_y}{2k^2}$ and $\xi_{21k}^x=-\frac{k_y}{2k^2}$ which are simply stating a sign change for traversing from one band to another. From Eq.(\ref{formula_conductivity}) the interband linear conductivity expression is written as,
\begin{equation}
\sigma_{inter}^{\!(1)xx}(\omega)=-e^2\int{\frac{d\textbf{k}}{4\pi^2}v_{12k}^x\mathcal{S}_{12k}^{ \!(1);x}(\omega)}-e^2\int{\frac{d\textbf{k}}{4\pi^2}v_{21k}^x\mathcal{S}_{21k}^{ \!(1);x}(\omega)}
\end{equation}
The interband frequency coefficient values are now placed in the above expression.
\begin{multline}
\begin{split}
\sigma_{inter}^{\!(1)xx}(\omega) = -e^2\int{\frac{d\textbf{k}}{4\pi^2}} \bigg[\frac{-v_{12k}^x\xi_{12k}^x\left(f_{1k}-f_{2k}\right)}{(\hbar\omega+i\Gamma_e^{\!(1)}-2\varepsilon)} \\
+\frac{v_{21k}^x\xi_{21k}^x\left(f_{1k}-f_{2k}\right)}{(\hbar\omega+i\Gamma_e^{\!(1)}+2\varepsilon)}\bigg]
\end{split}
\label{B8}
\end{multline}
If the interband velocity  matrix elements are substituted in the numerator with some simple steps the expression can be rewritten as,
\begin{multline}
\begin{split}
\sigma_{inter}^{\!(1)xx}(\omega) = \frac{ie^2}{2}\int{\frac{d\textbf{k}}{4\pi^2} \left(\frac{k_y}{k}\right)^2 \frac{v_F}{k}\frac{\left(f(+\varepsilon)-f(-\varepsilon)\right) }{(\hbar\omega+i\Gamma_i^{\!(1)})}} \\
\times \left[\frac{1}{(\hbar\omega+i\Gamma_e^{\!(1)}-2\varepsilon)}+\frac{1}{(\hbar\omega+i\Gamma_e^{\!(1)}+2\varepsilon)}\right]
\end{split}
\label{B9}
\end{multline}
Again, the transformation from momentum space to energy space is carried out to produce the other half of the result obtained by Kubo expression,
\begin{multline}
\sigma_{inter}^{\!(1)xx}(\omega) = \frac{ie^2}{\hbar\pi}\int\limits_{0}^{\infty}{\frac{d\varepsilon}{2}\left(f(+\varepsilon)-f(-\varepsilon)\right)} \\
 \left[\frac{1}{(\hbar\omega+i\Gamma_e^{\!(1)}-2\varepsilon)}+\frac{1}{(\hbar\omega+i\Gamma_e^{\!(1)}+2\varepsilon)}\right]
\label{Kubo_inter}
\end{multline}
By merging the Kubo expressions from above, one can evaluate the complete linear response in the presence of scattering. This approach has been extensively documented in several reputable manuscripts (\cite{CSV15}, \cite{hanson}). The equivalence between the solution obtained and the expression without accounting for any scattering phenomena \cite{CSV141} can be formally demonstrated through some easy mathematical steps. Another important criteria which needs to be considered during vanishing of the interband phenomenological constant towards the positive side otherwise the entire has to be taken in the reverse order which is not the target of our study. Furthermore, it should be noted that all the calculations hold true in the vicinity of absolute zero Kelvin. However, if necessary, appropriate adjustments to the chemical potential can be made to accommodate variations outside this temperature range. The indirect dependence of the scattering constant on temperature can be elucidated by expressing it in terms of curve fitting. Employing this approach, a more lucid comprehension of the thermal effects is attained, enabling the observation of temperature-related patterns and behaviors in the scattering constant.


%%%%%%%%%%%%%%   Bibliography   %%%%%%%%%%%%%%
\normalsize
\bibliography{preprint}

%%%%%%%%%%%%  Supplementary Figures  %%%%%%%%%%%%
%\clearpage

%%%%%%%%%%%%%%%%   End   %%%%%%%%%%%%%%%%
%\end{multicols}  % Method B for two-column formatting (doesn't play well with line numbers), comment out if using method A
\end{document}
